\scanpage{319}
\mspage{269}

Ce ``déploiement'' est un diagramme \inserted{infini} dont toutes les suites infinies, dans les trois directions des axes \unsure{de|des} coordonnées, sont des ``suites de suspension''. Je représente la portion utile du déploiement

\diagram{\(\begin{tikzcd}[ampersand replacement=\&,row sep=tiny,column sep=tiny]
{}\&{}\&{}\&{}\&{}\&{}\&\unsure{U|u}\arrow[dd,"{\mu}"]\\
{}\&X_{\deleted{\unsure{3}}2}\arrow[ur]\arrow[rr]\arrow[dd,dashed]\&{}\&X_3\arrow[ur]\arrow[rr,dash]\arrow[dd,dashed]\&{}\&T\arrow[ur]\arrow[dd]\&{}\\
X_0\arrow[ur]\arrow[rr]\arrow[dd]\&{}\&X_1\arrow[ur,dash]\arrow[rr]\arrow[dd]\&{}\&S\arrow[ur,"{\lambda}"]\arrow[dd]\&{}\&U'\arrow[dd]\\
{}\&Y_2\arrow[ur,shorten >=5pt]\arrow[rr,dashed]\arrow[dd,dashed]\&{}\&Y_3\arrow[ur,dashed,shorten >=6pt]\arrow[rr,dashed]\arrow[dd,dashed]\&{}\&\deleted{\unsure{\overline T}}T'\arrow[ur]\arrow[dd]\&{}\\
Y_0\arrow[ur,dash]\arrow[rr]\arrow[dd]\&{}\&Y_1\arrow[ur]\arrow[rr]\arrow[dd]\&{}\&S^{\unsure{\prime}}\arrow[ur,"{\lambda'}"]\arrow[dd]\&{}\&U''\\
{}\&Z_2\arrow[ur,shorten >=5pt]\arrow[rr,dashed]\&{}\&Z_3\arrow[dl]\arrow[ur,dashed,shorten >=6pt]\arrow[rr,dashed]\&{}\&T''\arrow[ur,line width=1pt]\&{}\\
Z_0\arrow[ur]\arrow[rr]\&{}\&Z_1\arrow[rr,dash]\&{}\&S''\arrow[ur,dash,"{\lambda''}"]\&{}\&{}
\end{tikzcd}\)}

La condition que $X$ soit $W$-cocartésien s'exprime par $\lambda\in W$, ou encore par $U\deleted{\unsure{\in}}\in\mathrm{Asph}_W$ (car $U$ s'identifie à \deleted{\unsure{$Z$|la}} \unsure{cofibre} homotopique $Z(\lambda)$ de $\lambda$ \unsure{\ldots}), de $=$ $Y$ $W$-cocartésien s'exprime par $U'\in\mathrm{Asph}_W$, et $Z$ $W$-cocartésien par $U''\in\mathrm{Asph}_W$. Or $U''\sim Z(\mu)$, où $\mu:U\longrightarrow U'$ est dans $W$ si $U,U'$ sont $W$-asphériques, donc $U''=Z(\mu)$ est alors $W$-asph. Cela prouve le lemme.

Cela montre que, moyennant $W(7\ \mathrm{bis})$ (en fait, on a utilisé seulement la partie ``universelle'' de $W(7\ \mathrm{bis})$, savoir que dans

\diagram{\(\deleted{\unsure{\#}}\quad\begin{tikzcd}[ampersand replacement=\&,row sep=tiny,column sep=small]
T_0\arrow[r]\arrow[d,"i"']\&T_2\arrow[d,"{i'}"]\\
T_{\unsure{1}}\arrow[r]\&T_{\unsure{3}}
\end{tikzcd}\)}

diagr. cartésien d'imm. ouvertes, $i\in W\Longrightarrow i'\in W$ (au lieu de $\Longleftrightarrow$),
% archive sheet number: 319
